% ===== §9 The Feminist Perspective =====
\section{The Feminist Perspective: Relational Subjectivity, Care Ethics, and the Non-Phallic Wealth}
\label{p15:sec:feminism}

\noindent
The feminist tradition offers perspectives on wealth in intimate relations that neither political economy nor psychoanalysis fully captures. Where political economy asks who produces and who owns, and where psychoanalysis asks what subjective structure the relation to wealth involves, feminist theory asks: whose labour has been systematically made invisible, whose subjectivity has been structurally denied, and what alternative forms of subjectivity and wealth might be possible if the dominant framework is challenged?

This section develops the feminist perspective in five movements. It begins with the structural problem of subjectivity in intimate relational life, showing how the dominant framework has denied genuine subjectivity to those who perform the labour of generating relational value. It then examines three feminist attempts to solve this problem---the first-wave liberal strategy, the second-wave critique of the liberal solution, and the care ethics alternative---showing how each contributes to and falls short of the GRW framework's account of relational subjectivity. Finally, it develops the feminist implications of GRW theory itself, arguing that GRW is constitutively a non-phallic wealth form that both requires and enables a genuinely relational subjectivity---and that this has implications not only for intimate relations but for the broader political economy of social reproduction.

\subsection{The Structural Problem: Whose Subjectivity?}
\label{p15:subsec:subjectivity_problem}

\noindent
The dominant frameworks for understanding intimate relations---from the legal and economic frameworks that organise their material dimensions to the philosophical and psychoanalytic frameworks that theorise their subjective dimensions---share a structural feature that feminist theory has consistently identified and challenged: they take the \emb{individual subject} as their fundamental unit, and they implicitly model this individual subject on the figure of the autonomous, self-determining, economically independent agent---a figure that has historically been identified with the masculine position in the social division of labour.

This structural feature has a specific consequence for the understanding of wealth in intimate relations: the wealth generated by those who perform the relational labour of care, emotional support, and attentive co-presence---labour historically performed predominantly by women---is systematically invisible in frameworks organised around the individual subject's accumulation, exchange, and enjoyment of wealth. Federici's analysis of social reproduction \parencite{federici2004} named this invisibility and located it in the structural requirements of capitalist production: capitalism requires the reproduction of labour power, which requires reproductive and care labour, which must therefore be performed without being recognised as wealth production, because its recognition would require its compensation and thereby undermine the extraction of surplus value from the workers whose reproduction it supports.

But the feminist critique goes deeper than the political economic level. The invisibility of care labour as wealth production is not merely an economic injustice; it is an \emb{ontological injustice}---a misrecognition of what the care labourer is and what they do. The woman whose emotional labour, attentive co-presence, and relational maintenance generates the GRW of the shared life is not merely economically exploited; she is ontologically misrecognised: her activity is not seen as productive, her contributions are not counted as wealth, and her subjectivity---constituted through and in the activity of relational care---is not recognised as a genuine form of human subjectivity.

\subsection{The First-Wave Liberal Strategy and Its Limits}
\label{p15:subsec:firstwave}

\noindent
The first-wave feminist response to this structural problem was to demand for women the same individual subjectivity that men had been granted: the right to own property, to enter contracts, to vote, to participate in public life as autonomous, self-determining agents. This strategy was politically crucial and historically transformative: it secured for women legal and economic protections that are genuine conditions for a minimally adequate intimate relational life.

But the first-wave strategy has a structural limitation that the GRW framework makes visible: it sought to secure for women the same \emb{individual} subjectivity that men possessed, thereby accepting the dominant framework's individualist model of subjectivity as the standard to be met. Women were to be recognised as individual subjects in the same sense that men were individual subjects---autonomous, self-determining, economically independent---rather than challenging the model of subjectivity itself.

This acceptance of the individualist model has a specific consequence for the understanding of wealth in intimate relations: it leaves intact the framework in which care labour and relational maintenance are not recognised as wealth production. If the standard of subjectivity is individual autonomy and self-determination, then the activity of relational care---which is constitutively not autonomous (it is responsive to the other) and not self-determining (it is shaped by the needs and dynamics of the relational field)---will continue to appear as a lesser form of activity, a deviation from the standard rather than a genuine alternative to it.

\subsection{The Second-Wave Critique and the Problem of the Alternative}
\label{p15:subsec:secondwave}

\noindent
The second-wave feminist tradition recognised the limitation of the first-wave strategy and developed a more radical critique of the dominant model of subjectivity. For Simone de Beauvoir, the fundamental problem was not the exclusion of women from a form of subjectivity they were otherwise capable of exercising, but the construction of women as \emb{Other} within a framework in which the subject---the position of genuine agency, freedom, and self-determination---was constitutively masculine \parencite{held2006}. Liberation, on this account, required not just inclusion in the existing subject position but the transformation of subjectivity itself: the recognition of women as genuine subjects, which required challenging the model of subjectivity that had been constructed in opposition to the feminine.

But the second-wave critique faced a problem of the alternative: if the dominant model of subjectivity is to be challenged, what is to replace it? The second-wave tradition largely answered this question by proposing a revised model of individual subjectivity---one that was more embodied, more emotionally engaged, more relational than the Cartesian model---but that retained the individual as its fundamental unit. Even in its most relational forms, second-wave feminism tended to understand the self as an individual self that is constituted through relations, rather than as a relational field that generates selves through co-evolutionary dynamics.

The GRB framework proposes a more radical answer to the problem of the alternative: not a revised model of individual subjectivity but a \emb{relational ontology} in which the relational field is the primary unit and individual subjectivity is a product of relational dynamics. This is not merely a philosophical position but a political one: it implies that the recognition of care labour as wealth production requires not only the redistribution of economic resources but the transformation of the ontological framework within which wealth and production are understood.

\subsection{Care Ethics and the Relational Alternative}
\label{p15:subsec:care_ethics}

\noindent
Carol Gilligan's groundbreaking study of women's moral reasoning \parencite{gilligan1982} identified an alternative to the dominant model of moral subjectivity that is directly relevant to the GRW framework. Where the dominant model (Kohlberg's stages of moral development) understood moral maturity as the capacity for abstract, universal, rule-governed reasoning---the capacity to apply general principles to particular cases without being distorted by particular attachments---Gilligan found that many of her female respondents reasoned in a fundamentally different way: not from general principles to particular cases, but from the particular relationships and particular responsibilities that constituted their lives to the responses that those relationships called for.

Gilligan called this the \emb{ethics of care}: a form of moral reasoning that begins not with universal principles but with particular relationships, that understands moral agency not as autonomy from attachment but as responsiveness to the demands of specific relationships, and that values not the impartiality of the universal moral law but the attentiveness and responsiveness that particular relationships require.

The ethics of care is not merely a different style of moral reasoning; it implies a different model of subjectivity. The care-ethical subject is not the autonomous individual of liberal theory, constituted independently of their relationships and choosing to enter them; it is a \emb{relational subject}, constituted through and in their particular relationships, whose agency is the agency of responsive attentiveness rather than autonomous self-determination.

Nel Noddings's development of care ethics \parencite{noddings1984} makes the relational subjectivity of the care-ethical framework even more explicit. For Noddings, caring is not a one-way relation of the carer toward the cared-for but a \emb{two-way relational dynamic}: genuine caring requires the carer to receive the cared-for's response and to allow that response to modify the caring activity. The carer who cares without receiving---who provides care without attending to the cared-for's reception of it---is not genuinely caring in Noddings's sense; they are performing care, which is a different and lesser thing. Genuine caring is a form of co-evolutionary engagement: the carer and the cared-for are coupled in a dynamic in which each modifies the other's activity, producing a caring relationship that is irreducible to either party's individual contribution.

Virginia Held's political extension of care ethics \parencite{held2006} argues that care is not merely a private relational practice but a political principle: that the relations of care that sustain human life should be the model for political arrangements, rather than the market relations of self-interested exchange or the contractual relations of autonomous agents. This political extension is directly relevant to the GRW framework: it suggests that the transformation of intimate relational life along GRW lines is not merely a personal project but a political one, requiring the transformation of the social and economic structures that currently prevent care from being recognised and sustained as the primary form of human value production.

\begin{claim}[Care ethics and GRW: the convergence]
The ethics of care tradition---Gilligan's identification of care-ethical reasoning, Noddings's account of caring as a two-way relational dynamic, and Held's political extension of care ethics---provides an independent philosophical grounding for the GRW framework's account of relational subjectivity and relational wealth. Care labour is GRW production; the care-ethical subject is the relational subject of GRW generation; and the political transformation that care ethics calls for---the recognition of care as the primary form of human value production---is the political dimension of the GRW framework's theoretical project.
\end{claim}

\subsection{Relational Subjectivity: The Third Way}
\label{p15:subsec:relational_subjectivity}

\noindent
The GRW framework proposes what we call \emb{relational subjectivity} as an alternative to both the liberal individual subject (the autonomous, self-determining agent of the dominant framework) and the care-ethical subject (the responsive, attentive caregiver of Gilligan and Noddings). Relational subjectivity is not a compromise between these two positions but a genuinely different ontological proposal: the claim that the primary unit of subjectivity is not the individual but the \emb{relational field}, and that individual subjects are products of relational dynamics rather than their preconditions.

This proposal has several implications for the understanding of gender and wealth in intimate relations.

\emb{First}, relational subjectivity dissolves the opposition between the autonomous individual and the relational caregiver that has structured feminist debate since de Beauvoir. If the primary subject is the relational field, then neither the autonomous individual nor the responsive caregiver is the fundamental form of subjectivity; both are specific modes of participation in relational dynamics, and the question is not which is more genuinely human but which modes of participation are most generative of GRW.

\emb{Second}, relational subjectivity reconceives the problem of care labour's invisibility. On the individualist framework, care labour is invisible because it does not produce individual wealth---it does not accumulate in the individual's account, does not build individual capital, does not enhance individual autonomy. On the relational subjectivity framework, care labour is the primary form of wealth production precisely because it is the activity that most directly generates and sustains the relational field's GRW. The invisibility of care labour is therefore not a reflection of its unproductiveness but of the dominant framework's inability to see what it produces---a form of epistemic injustice that \xref{p15:sec:ethics} analyses in detail.

\emb{Third}, relational subjectivity provides the ontological basis for a non-essentialist feminist politics of intimate life. The claim that care labour generates GRW does not require the claim that women are naturally more caring than men, or that care is essentially feminine. It requires only the empirical observation that care labour has historically been performed predominantly by women, and the normative claim that the distribution of care labour should be organised to maximise the relational field's GRW generation rather than to reproduce existing gender hierarchies. A genuinely relational subjectivity is gender-neutral in its structure---both parties to the relational field are constituted through it, both contribute to its dynamics, and both benefit from its GRW. The gendered distribution of care labour is a historical contingency that the GRW framework, by making care labour visible as the primary form of wealth production, provides resources to challenge.

\subsection{The Non-Phallic Wealth and Its Political Implications}
\label{p15:subsec:nonphallic}

\noindent
GRW is constitutively non-phallic: it cannot be had or lacked, it is not organised by the logic of symbolic power, and it does not promise the completeness that phallic objects promise. This non-phallic character has implications that extend beyond the intimate relational field into the broader political economy of social life.

The phallic economy---the economy organised around the logic of having and lacking, in which wealth is the medium of symbolic power and its accumulation is the pursuit of phallic completeness---is not gender-neutral. The phallic position has been historically identified with the masculine, and the wealth forms that most directly instantiate phallic logic---gold, money, property, productive capital---have been historically controlled predominantly by men. The feminist political economy of intimate relations is therefore inseparable from the critique of phallic wealth: challenging the organisation of intimate relations around the logic of having and lacking is inseparable from challenging the gendered distribution of symbolic power in the broader social field.

GRW's non-phallic character makes it a potentially \emb{liberatory} form of wealth---not because it eliminates the need for material wealth (which it does not) but because it provides a form of wealth that is not organised by the phallic logic of symbolic power. A relational field rich in GRW is not a field in which one party has more symbolic power than the other; it is a field in which both parties participate in the co-evolutionary dynamics that generate a wealth that belongs to neither of them individually. This is the economic correlate of what the care ethics tradition calls genuine caring: not the unilateral provision of care by one party for another, but the co-evolutionary dynamic in which both parties are simultaneously carers and cared-for, both generators of GRW and beneficiaries of it.

The political implication is significant: if GRW is the primary form of wealth in intimate relations, and if GRW is non-phallic in its structure, then the organisation of intimate relations around GRW rather than around material wealth would represent a transformation of the gendered power structure of intimate life that goes beyond what either liberal feminism (which seeks equal access to phallic wealth) or radical feminism (which critiques the phallic structure itself) has fully articulated. It is the feminist political economy of the \emph{between}: not the redistribution of phallic wealth between individual subjects, but the cultivation of a form of wealth that belongs to the relational field and that is, by its very structure, immune to the phallic logic of having and lacking.

Sara Ruddick's account of maternal thinking \parencite{ruddick1989} provides a further dimension of this feminist political economy. Ruddick argues that the practice of caring for children---attending to their particular needs, responding to their particular vulnerabilities, fostering their growth without controlling its direction---develops a specific form of practical reasoning that she calls \emb{attentive love}: a form of knowledge that is constitutively relational, responsive, and non-possessive. Attentive love is, in Ruddick's account, a politically significant epistemic practice: it is a form of knowing that the dominant framework cannot recognise as knowledge, because it does not fit the model of detached, universal, rule-governed cognition that the dominant framework treats as the standard.

In the GRW framework, attentive love is not merely an epistemic practice but a \emb{wealth generation practice}: the practice of attending to the particular needs and dynamics of the relational field, responding to its contingent events with genuine co-present receptivity, and fostering its growth without claiming its products---the practice of \emph{xuande} in its most fully developed form. The feminist insight that this practice has been undervalued and made invisible is the feminist political economy of GRW: the claim that the most important form of wealth production in human social life has been rendered invisible by the dominant framework's inability to see beyond the phallic logic of having and lacking.
